\section{Loop-space recovery}
\label{sec:loops}

\subsection{Loops as equivariant automorphisms}
At the regular basepoint, the inverse of \eqref{eq:sip-torsors} gives
\begin{equation}
 \Loop(G\mathrm{Tors}_{L,\ell},b_G)\simeq\operatorname{Aut}_G(G),
 \label{eq:loops-aut}
\end{equation}
where $G$ has the regular left action.  In the source this is
\agda{torsorLoopEquivAut}, obtained by applying \agda{invEquiv} to
\agda{torsorPathEquiv}.
\formal{\texttt{src/TorsorDelooping/Delooping.agda}}
{\texttt{torsorLoopEquivAut}}
{The equivalence is from loops to structured equivariant self-equivalences of the regular torsor.}

\subsection{Evaluation and right translations}
Evaluation at the unit is the map $E(f)=f(1)$ from equivariant automorphisms
to $G$.  Its inverse sends $h:G$ to the right translation
\begin{equation}
 R_h(x)=xh.
 \label{eq:right-translation}
\end{equation}
It is equivariant for the left action because $R_h(gx)=(gx)h=g(xh)$.  It is
invertible through $R_{h^{-1}}$.  Conversely, equivariance forces
\begin{equation}
 f(x)=f(x1)=xf(1)=R_{f(1)}(x).
 \label{eq:equivariant-determined}
\end{equation}
Function extensionality proves equality of the underlying equivalences, and
\agda{Σ≡Prop} lifts it through the propositional equivariance field.  Thus
evaluation and right translation form an isomorphism.
\formal{\texttt{src/TorsorDelooping/Delooping.agda}}
{\texttt{regularAutToGroup}; \texttt{rightRegularAut}; \texttt{regularAutUnderlying}; \texttt{regularAutRet}; \texttt{regularAutIsoGroup}}
{These declarations implement evaluation, right translation, extensional reconstruction, and the resulting \texttt{Iso}.}

Composing \eqref{eq:loops-aut} with this isomorphism yields
\begin{equation}
 \Loop(G\mathrm{Tors}_{L,\ell},b_G)\simeq G.
 \label{eq:generic-loop}
\end{equation}
\formal{\texttt{src/TorsorDelooping/Delooping.agda}}
{\texttt{torsorLoopEquivGroup}}
{Its implementation is \texttt{compEquiv torsorLoopEquivAut (isoToEquiv regularAutIsoGroup)}.}

\subsection{Order conventions and the formal boundary}
There is no hidden opposite group in this calculation.  If composition means
first $R_g$ and then $R_h$, then
\begin{equation}
 R_h(R_g(x))=(xg)h=x(gh)=R_{gh}(x).
\end{equation}
This is the correct order for right translations acting as equivariant maps
of a left torsor.  The generic CQTS development has a separate regular-loop
coherence used for its $\mathsf{EM}_1$ construction, but
\agda{torsorLoopEquivGroup} is not packaged as a group-isomorphism theorem.

\begin{warning}[No group-level loop-recovery claim]
Loop concatenation gives the source of \eqref{eq:generic-loop} a group
structure, but the formal endpoint supplies only an equivalence of underlying
types.  No formal theorem here asserts preservation of multiplication, unit,
or inverse.
\end{warning}

For $G=\CHunits$, \eqref{eq:generic-loop} is the local declaration
\begin{equation}
 \Loop(\BG\CHunits,b_{\CHunits})\simeq\CHunits.
\end{equation}
\formal{\texttt{src/ConstructiveComplex/Delooping.agda}}
{\texttt{complexUnitsLoopEquiv}}
{This is definitionally \texttt{torsorLoopEquivGroup ComplexUnitsGroup}; its codomain is \texttt{ComplexUnits}.}
